% =====================================================================
%  Báo cáo ngắn: kết quả B0 (MoSS) và M1 -- M3 trên POJ-104
%  Biên dịch bằng XELATEX (cần font hỗ trợ tiếng Việt):
%      xelatex -interaction=nonstopmode 012_m1_m2_m3_summary.tex
%  Chạy 2 lần để chốt số trang.
%
%  Số liệu lấy từ:
%    - results/*/metrics.json              (P/R/F1/ROC-AUC trên test)
%    - results/uncertainty/uncertainty_report.json  (bootstrap, McNemar)
%  Không gõ tay số liệu vào bảng.
% =====================================================================
\documentclass[11pt,a4paper]{article}

\usepackage{fontspec}
\setmainfont{Times New Roman}
\setsansfont{Arial}
\newfontfamily\monofont{Consolas}

\usepackage[margin=2.3cm]{geometry}
\usepackage{booktabs}
\usepackage[table]{xcolor}
\usepackage{enumitem}

\definecolor{headblue}{RGB}{30,60,110}

\setlength{\parskip}{5pt}
\setlength{\parindent}{0pt}

\newcommand{\hd}[1]{\textcolor{headblue}{\textbf{#1}}}

\title{\vspace{-2.0cm}\textbf{\color{headblue}Kết quả so sánh trên dataset POJ-104}\\[6pt]
       \large Baseline MOSS và ba mô hình M1 -- M3}
\author{Nhóm nghiên cứu}
\date{}

\begin{document}
\maketitle
\vspace{-1.2cm}

\section{Việc cần làm và dữ liệu}

Bài toán: cho hai hàm code, có phải là bản sao của nhau hay không. Dùng dataset
POJ-104, gồm 51.991 hàm C++. Ghép thành 66.000 cặp, mỗi cặp có nhãn ``bản sao''
hay ``không bản sao''. Chia ba phần: train 45.000 cặp, validation 10.000, test
11.000. Tập test cân bằng, có 5.500 cặp dương và 5.500 cặp âm.

Ngưỡng quyết định của mọi mô hình được chọn trên \emph{validation} rồi giữ
nguyên, không tinh chỉnh lại trên test.

\section{Các mô hình}

\begin{itemize}[leftmargin=1.3em, itemsep=2pt]
    \item \textbf{B0 -- MOSS.} Thuật toán Winnowing của Schleimer, Wilkerson \& Aiken
          (2003), cài lại chạy cục bộ (EXP-013). Đây là baseline tham chiếu.
    \item \textbf{M1 -- CodeBERT.} Mô hình ngôn ngữ cheo token, fine-tune để so hai hàm.
    \item \textbf{M2 -- GraphCodeBERT.} Giống M1 nhưng có thêm thông tin cấu trúc
          (quan hệ cây cú pháp, dữ liệu luồng) nối các token.
    \item \textbf{M3 -- ASTNN.} Rút cây cú pháp thành đồ thị, đưa qua mạng nơ-ron
          đệ quy để tạo vector đặc trưng. Không dùng token thô.
\end{itemize}

\section{Kết quả trên tập test}

Cột F1 ghi dạng \texttt{mean $\pm$ SD}. Ý nghĩa: lấy tập test hiện có, bốc
ngẫu nhiên lại 2.000 lần, tính F1 mỗi lần, rồi lấy trung bình và độ lệch chuẩn.

\begin{center}
\begin{tabular}{lcccc}
\toprule
\hd{Mô hình} & \hd{Precision} & \hd{Recall} & \hd{F1} & \hd{ROC-AUC} \\
\midrule
B0 -- MOSS              & 0.646 & 0.722 & $0.682 \pm 0.005$ & 0.731 \\
M1 -- CodeBERT          & 0.931 & 0.890 & $0.910 \pm 0.003$ & 0.973 \\
M2 -- GraphCodeBERT     & 0.883 & 0.933 & $0.907 \pm 0.003$ & 0.971 \\
M3 -- ASTNN             & 0.641 & 0.903 & $0.750 \pm 0.004$ & 0.795 \\
\bottomrule
\end{tabular}
\end{center}

Vài điều đọc được từ bảng:

\begin{itemize}[leftmargin=1.3em, itemsep=2pt]
    \item M1 và M2 gần như ngang nhau. Chênh lệch F1 chỉ 0.003, nhỏ hơn cả độ lệch
          chuẩn. Thêm thông tin cấu trúc của M2 không giúp ở bài toán này.
    \item M1 thiên về precision, M2 thiên về recall. M1 bỏ sót nhiều cặp bản sao hơn,
          M2 đoán nhầm nhiều cặp không phải bản sao hơn.
    \item MOSS kém hơn M1 khá xa (F1 0.682 so với 0.910). Lý do dễ hiểu: MOSS so
          các đoạn chữ giống hệt từng ký tự, mà trong POJ-104 hai hàm là bản sao
          thì tên biến và cấu trúc đã bị đổi khá nhiều.
    \item M3 nhớ nhiều bản sao (recall 0.903) nhưng hay đoán nhầm (precision 0.641).
\end{itemize}

\section{M1 và M2 có thật sự khác nhau không?}

Không thể trả lời bằng cách nhìn hai cột $\pm$ SD rồi kết luận, vì SD ở đây đo
độ lớn của tập test chứ không đo sai khác giữa hai mô hình. Cách đúng là so hai
mô hình trên \emph{cùng một} tập test.

\begin{center}
\begin{tabular}{lccc}
\toprule
\hd{Cặp so sánh} & \hd{$\Delta$F1} & \hd{Khoảng tin cậy 95\%} & \hd{McNemar} \\
\midrule
M1 -- M2  & $-0.003$ & $[-0.0085,\ +0.0028]$ & $p = 0.013$ \\
M1 -- B0  & $+0.228$ & $[+0.2175,\ +0.2388]$ & $p < 10^{-300}$ \\
M1 -- M3  & $+0.160$ & $[+0.1507,\ +0.1688]$ & $p < 10^{-300}$ \\
\bottomrule
\end{tabular}
\end{center}

Hai cột cuối lại không thống nhất, và điều đó có nghĩa:

\begin{itemize}[leftmargin=1.3em, itemsep=2pt]
    \item Khoảng tin cậy của $\Delta$F1 \emph{chứa} 0, nên không đủ bằng chứng để nói
          M1 khác M2 về F1.
    \item McNemar lại ra $p = 0.013$: hai mô hình \emph{có} khác nhau về số cặp dự
          đoán đúng.
    \item Hai kết quả này không mâu thuẫn. Chúng khác nhau vì đo hai thứ khác nhau:
          McNemar hỏi ``tổng số dự đoán đúng có khác không'', còn khoảng tin cậy hỏi
          ``F1 có khác không''. M1 và M2 sai ở \emph{hướng khác nhau} (M1 hơi thiên
          precision, M2 thiên recall) nên trong F1 hai sai lệch đó gần như triệt tiêu.
\end{itemize}

Kết luận đúng nhất là: \textbf{M1 và M2 không phân biệt được về F1}. Chọn cái nào
cũng được về mặt F1; nếu cần recall cao thì chọn M2.

Với MOSS và M3 thì khoảng tin cậy không chứa 0, tức là khác biệt thật, không phải
nhiễu.

\section{Hạn chế cần nói rõ}

\begin{itemize}[leftmargin=1.3em, itemsep=2pt]
    \item Mỗi mô hình mới chạy \emph{một lần}, nên chưa có độ lệch chuẩn do biến
          thiên giữa các lần train. Muốn có con số đó thì phải train lại vài seed;
          với M1/M2 mỗi lần mất khoảng 7 giờ GPU. Hiện tại $\pm$ SD trong bảng chỉ
          phản ánh độ lớn của tập test.
    \item M2 vẫn đang train dở (3 epoch, F1 trên validation còn đang tăng), nên số
          0.907 này có thể còn tăng nữa.
    \item M1 và M2 bị cắt còn 256 token. Hàm POJ-104 trung bình dài hơn thế, nên
          phần đuôi hàm bị mất.
    \item M3 có thể chưa train tới đủ. Số của M3 nên xem là tạm.
\end{itemize}

\section{Đang làm tiếp}

M4 (CPG-GNN) đang xây dựng: cần dựng code property graph bằng Joern trên Kaggle
trước, nên hiện chưa có kết quả. Sau khi M1 -- M3 chốt xong sẽ cho M4 chạy và so
thêm vào bảng này.

\end{document}